\scanpage{12}
% archive sheet number: 11

\uline{Terminologie} Variété d'appui d'un convexe $C$ $=$ variété affine de la forme $E_F$. Les variétés d'appui sont en corr. $\unsure{1\text{-}2|1\text{-}1}$ avec les facettes, relation d'ordre la relation d'inclusion

\margin{\begin{tikzpicture}[x=1em,y=1em]\draw (.15,6.4)--(.16,5.5); \draw (.07,5.54) .. controls (0,4.3) and (-.03,1.3) .. (.18,0);\end{tikzpicture}} \uline{Th} Soit $C,C'$ deux secteurs polyédraux polaires l'un de l'autre de $E,E'$ resp$^{\mathrm t}$. Alors il existe une unique \deleted{\unsure{corresp}} application $\unsure{\varphi|\phi}$ de l'ens des facettes $F$ de $C$ telles que $0\notin\overline F$ (i.e. $0\in E_F$) vers l'ens des facettes $F'$ de $C'$ telles que $0\notin\overline{F'}$ (i.e. $0'\in E_{F'}$), caractérisée par la propriété que \deleted{si $F'=\varphi(F)$,} \deleted{\unsure{on a} $E$}

\[
 E_{\varphi(F)=F'}=(E_F)^{\unsure{\vee|\circ|0}}
\]

Cette application est bijective, et renverse la relation d'ordre entre facettes

\uline{Cor. 1} L'ens des facettes d'un secteur polyédral est fini (prendre l'origine \uline{intérieure} au secteur polyédral, et appliquer dualité)

\uline{Cor 2} Supposons $C,C'$ cpct \deleted{\unsure{resp.}} i.e. \deleted{$\unsure{\in|C|E}$} $0_E,0_{E'}$ \unsure{intérieurs} à $C,C'$ resp$^{\mathrm t}$, i.e. $C,C'$ des polyèdres. Alors $\varphi:\unsure{\operatorname{Facl}|\operatorname{Fac}}(C)\simeq\unsure{\operatorname{Facl}|\operatorname{Fac}}(C')$
